\documentclass[11pt,a4paper]{article}

\usepackage[utf8]{inputenc}
\usepackage[T1]{fontenc}
\usepackage[spanish,provide=*]{babel}
\usepackage[margin=2.5cm]{geometry}
\usepackage{graphicx}
\usepackage{xcolor}
\usepackage{listings}
\usepackage{tikz}
\usepackage{hyperref}
\usepackage{fancyhdr}
\usepackage{enumitem}
\usepackage{longtable}
\usepackage{booktabs}
\usepackage{array}
\usetikzlibrary{arrows.meta, positioning, calc, shapes.geometric}

\definecolor{codebg}{RGB}{245,245,245}
\definecolor{codekeyword}{RGB}{0,0,150}
\definecolor{codestring}{RGB}{140,0,0}
\definecolor{codecomment}{RGB}{0,120,0}
\definecolor{uteqverde}{RGB}{0,90,60}
\definecolor{alerta}{RGB}{150,30,30}
\definecolor{okverde}{RGB}{0,110,60}

\lstdefinestyle{codigo}{
  backgroundcolor=\color{codebg},
  basicstyle=\ttfamily\footnotesize,
  keywordstyle=\color{codekeyword}\bfseries,
  stringstyle=\color{codestring},
  commentstyle=\color{codecomment}\itshape,
  numbers=left,
  numberstyle=\tiny\color{gray},
  breaklines=true,
  frame=single,
  rulecolor=\color{black!30},
  tabsize=2,
  showstringspaces=false,
  captionpos=b
}
\lstset{style=codigo}

\hypersetup{
  colorlinks=true,
  linkcolor=uteqverde,
  urlcolor=uteqverde,
  citecolor=uteqverde,
  pdftitle={SGED - Informe de la Practica Experimental Unidad IV},
  pdfauthor={Velez Lopez Ricardo Elias, Integrante A, Integrante B}
}

\pagestyle{fancy}
\fancyhf{}
\setlength{\headheight}{14pt}
\renewcommand{\headrulewidth}{0.4pt}
\lhead{SGED --- Pr\'actica Experimental Unidad IV}
\rhead{Aplicaciones Web --- PPA 2026-2027}
\cfoot{\thepage}

\newcommand{\ruta}[1]{\texttt{\footnotesize #1}}
\newcommand{\ok}{\textcolor{okverde}{\textbf{OK}}}
\newcommand{\pendiente}{\textcolor{alerta}{\textbf{pendiente}}}
\newcommand{\porintegrante}[1]{\textcolor{alerta}{\textbf{[Pendiente de #1]}}}

\begin{document}

% ============================================================
% PORTADA
% ============================================================
\begin{titlepage}
\centering
\vspace*{1cm}

{\LARGE\bfseries UNIVERSIDAD TÉCNICA ESTATAL DE QUEVEDO}\\[0.3cm]
{\large Facultad de Ciencias de la Ingeniería}\\[0.15cm]
{\large Carrera de Ingeniería en Software}\\[2cm]

{\Huge\bfseries\color{uteqverde} SGED}\\[0.4cm]
{\Large Sistema de Gestión para la\\Escuela Deportiva ProFútbol}\\[1.5cm]

{\LARGE\bfseries Práctica Experimental --- Unidad IV}\\[0.4cm]
{\large Informe Técnico Final}\\[0.3cm]
{\large Proyecto Fin de Curso --- Aplicaciones Web}\\[2cm]

\begin{tabular}{rl}
\textbf{Asignatura:} & Aplicaciones Web (Quinto nivel) \\[0.2cm]
\textbf{Docente:} & Dr. Gleiston Cicerón Guerrero Ulloa, Ph.D. \\[0.2cm]
\textbf{Periodo:} & PPA 2026-2027 \\[0.2cm]
\textbf{Etiqueta:} & \texttt{v1.0.0-rc} (candidata a versión estable) \\[0.2cm]
\textbf{\emph{Commit} del corte:} & \texttt{23deb19} \\[0.2cm]
\textbf{Repositorio:} & \url{https://github.com/Richiflo0o/SGED_APPWEB_PE_U4.git} \\
\end{tabular}\\[1.5cm]

{\large\bfseries Integrantes (grupo de la práctica):}\\[0.3cm]
{\large Velez Lopez Ricardo Elias}\\
{\large \porintegrante{Integrante A}}\\
{\large \porintegrante{Integrante B}}\\[1.5cm]

{\large Quevedo --- Los Ríos --- Ecuador --- 2026}

\vfill
\end{titlepage}

\tableofcontents
\newpage

% ============================================================
\section*{Resumen}
\addcontentsline{toc}{section}{Resumen}

Este informe documenta la Práctica Experimental de la Unidad IV del Proyecto
Fin de Curso (PFC): la verificación técnica del sistema \textbf{SGED}
(Sistema de Gestión para la Escuela Deportiva ProFútbol), aplicación web MVC
construida con Spring Boot 3.2.5, Angular y PostgreSQL 16. El trabajo se
organizó en bloques con autoría verificable por \emph{commit}: exposición de
la API REST documentada con Swagger/OpenAPI en \texttt{/api/docs} (Bloque 1),
anotaciones \texttt{@Tag/@Operation/@Schema} en controladores y DTOs (Bloque
2), auditoría de seguridad alineada con el OWASP Top 10 ampliada a A04, A06,
A07, A09 y XSS (Bloque 3), migraciones Flyway V1--V7 alineadas con el esquema
y el contenedor Docker reproducible (Bloque 4), consumo de una API externa de
datos de fútbol con caché Redis (Bloque 5.1) y la medición de rendimiento con
k6 (Bloque 5.3). La seguridad se evidencia con peticiones \texttt{curl}
reproducibles guardadas en \ruta{docs/mediciones/sec/}, incluidas las
cabeceras de seguridad, el bloqueo por fuerza bruta (429 tras seis intentos
fallidos) y la protección ante XSS con CSP y tipo de contenido
\texttt{application/json}. Este informe declara además, de forma explícita y
transparente, el uso de herramientas de asistencia por inteligencia
artificial en el desarrollo.

\section*{Abstract}
\addcontentsline{toc}{section}{Abstract}

This report documents the Experimental Practice of Unit IV of the Final
Course Project (PFC): the technical verification of the \textbf{SGED} system
(Sports Management System for the ProFútbol Sports School), a MVC web
application built with Spring Boot 3.2.5, Angular and PostgreSQL 16. The work
was organized in blocks with authorship verifiable by \emph{commit}: REST API
exposure documented with Swagger/OpenAPI at \texttt{/api/docs} (Block 1),
\texttt{@Tag/@Operation/@Schema} annotations in controllers and DTOs (Block
2), a security audit aligned with the OWASP Top 10 extended to A04, A06, A07,
A09 and XSS (Block 3), Flyway migrations V1--V7 aligned with the schema and a
reproducible Docker container (Block 4), external football-data.org API
consumption with Redis caching (Block 5.1), and load testing with k6 (Block
5.3). Security is evidenced with reproducible \texttt{curl} requests stored
in \ruta{docs/mediciones/sec/}, including security headers, brute-force
blocking (429 after six failed attempts) and XSS protection through CSP and
\texttt{application/json} content type. This report additionally declares,
explicitly and transparently, the use of artificial intelligence assistance
tools during development.

\newpage

% ============================================================
\section{Introducción}

La Práctica Experimental de la Unidad IV exige aplicar sobre el Proyecto Fin
de Curso las competencias de la asignatura Aplicaciones Web: una aplicación
web MVC completa con los módulos de la Entrega 1B funcionando, API REST
documentada con Swagger, consumo de una API externa con caché Redis,
seguridad según el OWASP Top 10, pruebas de carga y el informe técnico final.
Este repositorio (\url{https://github.com/Richiflo0o/SGED_APPWEB_PE_U4.git})
es un clon derivado de \texttt{DarwinSM21/SGED\_APPWEB} que conserva el
historial completo de las entregas 1A--3 y registra el trabajo de la práctica
en la rama \texttt{feature/u4-docs-seguridad}.

El trabajo se dividió en bloques. Los bloques 1--4 y la auditoría de
seguridad fueron desarrollados y verificados por Ricardo Velez Lopez
(autoría verificable en el historial de \texttt{git}); los bloques 5.1
(consumo de la API externa) y 5.3 (pruebas de carga) son responsabilidad de
los integrantes A y B respectivamente y quedan marcados en este informe con
los espacios reservados para sus resultados. La distribución de autoría se
resume en la sección~\ref{sec:autorias}.

\section{Metodología}

La práctica se desarrolló con verificación continua y evidencia
reproducible: cada bloque produce artefactos versionados en el repositorio
(evidencia cruda de \texttt{curl}, reportes de cobertura, salidas de k6), de
modo que el profesor puede reproducir cada resultado con los comandos del
\emph{Makefile} sin depender de la palabra de los integrantes.

\begin{center}
\begin{tabular}{@{}ll@{}}
\toprule
\textbf{Comando} & \textbf{Acción} \\
\midrule
\texttt{make up} & Levanta el sistema completo (Docker Compose) \\
\texttt{make test} & Pruebas JUnit 5 + cobertura JaCoCo \\
\texttt{make bench} & 3 corridas k6 (50 VUs, 30\,s) + análisis IC 95\,\% \\
\texttt{make audit} & Auditoría OWASP (A01--A07, A09, XSS) + SQL dinámico \\
\texttt{make clean} & Limpia contenedores, volúmenes y compilaciones \\
\bottomrule
\end{tabular}
\end{center}

Las credenciales semilla para probar el sistema son \texttt{admin} /
\texttt{Admin2026!} (definidas en \ruta{db/seed.sql}). El certificado TLS de
\texttt{https://localhost:8443} es autofirmado, generado solo para desarrollo
y evaluación.

\section{La API REST documentada (Bloques 1--2)}

SGED expone su API REST bajo el prefijo \texttt{/api}. La documentación se
publica con springdoc-openapi 2.3.0 en dos rutas:

\begin{itemize}[leftmargin=*]
  \item \textbf{Swagger UI}: \texttt{GET http://localhost:8080/api/docs} ---
    interfaz interactiva que permite explorar y probar cada \emph{endpoint}.
  \item \textbf{OpenAPI 3.0 en JSON}: \texttt{GET
    http://localhost:8080/api/api-docs} --- documento consumible por
    herramientas y generadores de clientes.
\end{itemize}

Ambas rutas están permitidas en \ruta{SecurityConfig.java} y accesibles sin
autenticación para su consulta. La configuración vive en
\ruta{application.yml}:

\begin{lstlisting}[language=Java,caption={Rutas de la documentaci\'on (Bloque 1)}]
springdoc:
  api-docs:
    path: /api/api-docs
  swagger-ui:
    path: /api/docs
\end{lstlisting}

\subsection{Anotaciones OpenAPI (Bloque 2)}

Cada controlador del dominio \texttt{academico}, \texttt{deportivo} y
\texttt{seguridad} está anotado con \texttt{@Tag} (agrupación), y cada
\emph{endpoint} con \texttt{@Operation} (descripción, resumen y respuestas
posibles). Los DTOs de entrada y salida (\texttt{EstudianteRequest},
\texttt{EstudianteResponse}, \texttt{LoginRequest}, \texttt{RegisterRequest},
\texttt{SesionResponse}, etc.) usan \texttt{@Schema} para documentar campos,
tipos y restricciones. Ejemplo:

\begin{lstlisting}[language=Java,caption={Anotaciones OpenAPI en EstudianteController}]
@Tag(name = "Estudiantes",
     description = "CRUD de estudiantes con paginacion, soft delete y conteos por categoria")
public class EstudianteController {

    @GetMapping
    @PreAuthorize("hasAnyRole('ADMINISTRADOR', 'ENTRENADOR', 'USER')")
    @Operation(summary = "Listar estudiantes paginados",
        description = "Devuelve una pagina de estudiantes ordenables.")
    @ApiResponses({
        @ApiResponse(responseCode = "200", description = "Pagina de estudiantes"),
        @ApiResponse(responseCode = "403", description = "Rol sin permiso")
    })
    public ResponseEntity<EstudiantePageResponse<EstudianteResponse>> listar(
        @RequestParam(defaultValue = "0") int page,
        @RequestParam(defaultValue = "10") int size,
        @RequestParam(defaultValue = "idEstudiante,asc") String[] sort) {
        ...
    }
}
\end{lstlisting}

La colección de Postman (\ruta{docs/postman/coleccion.json}) documenta el uso
manual de la API, y el atributo de calidad de la documentación se mapea en
\ruta{docs/iso25010-atributos-calidad.md} sobre el modelo ISO/IEC
25010~\cite{iso25010}.

% ============================================================
\section{Consumo de API externa con caché Redis (Bloque 5.1)}
\label{sec:apiexterna}
\porintegrante{Integrante A. Sección reservada.}

Esta sección documenta la integración de la API externa de datos de fútbol
(\texttt{football-data.org}) consumida por SGED, con la estrategia de caché
en Redis para limitar la cantidad de llamadas y el tiempo de respuesta.

\porintegrante{Integrante A: describir aquí el \emph{endpoint} consumido, el
recurso que resuelve en SGED, el TTL de 24\,h configurado en Redis, el
manejo de errores ante fallos de la API externa y la evidencia de una
petición real (salida de \texttt{curl} o captura).}

% ============================================================
\section{Seguridad: auditoría OWASP (Bloque 3)}
\label{sec:seguridad}

La seguridad se verificó con el \emph{script} \ruta{scripts/audit-owasp.sh},
que ejecuta peticiones \texttt{curl} reales contra el sistema levantado,
guarda la evidencia cruda en \ruta{docs/mediciones/sec/} con fecha y
\emph{commit}, y comprueba cada control en ambas direcciones: el
comportamiento seguro esperado \emph{y} que no se haya roto el uso legítimo.
Se cubren los diez riesgos del OWASP Top 10:2021~\cite{owasp2021} en su
alcance aplicable, más la categoría XSS (Cross-Site Scripting) de forma
explícita.

\subsection{A01 --- Broken Access Control}

Se comprueba con un usuario autenticado sin rol administrativo que las
operaciones administrativas devuelvan \texttt{403}, y que las lecturas
permitidas al rol \texttt{USER} sigan funcionando:

\begin{center}
\begin{tabular}{@{}ll@{}}
\toprule
\textbf{Petición} & \textbf{Resultado} \\
\midrule
\texttt{POST /api/estudiantes/operaciones/desactivar-categoria} & 403 \\
\texttt{GET /api/personas} & 403 \\
\texttt{GET /api/personas/cedula/0000000000} & 403 \\
\texttt{POST /api/categorias} & 403 \\
\texttt{DELETE /api/categorias/1} & 403 \\
\texttt{PUT /api/personas/1} & 403 \\
\texttt{GET /api/usuarios} & 403 \\
\midrule
\texttt{GET /api/categorias/activas} (uso legítimo) & 200 \\
\texttt{GET /api/estados\_generales} (uso legítimo) & 200 \\
\bottomrule
\end{tabular}
\end{center}

El error se devuelve como \texttt{application/problem+json} conforme al RFC
9457~\cite{rfc9457}, con \texttt{title: Forbidden} y el detalle «No tiene
permisos para acceder a este recurso». La auditoría limpia al inicio su
propio contador de intentos de login en Redis, porque de lo contrario el
propio control A07 bloquearía por IP la ejecución y A01 devolvería \texttt{401}
en todo, un resultado que parece correcto y no lo es (ver
sección~\ref{sec:reflexion}).

\subsection{A02 --- Cryptographic Failures}

La evidencia \ruta{docs/mediciones/sec/a02-tls.txt} muestra el apretón de
manos TLS contra el frontend: \textbf{TLS 1.3}, cifrado
\texttt{TLS\_AES\_256\_GCM\_SHA384}, intercambio de claves X25519 y firma
RSASSA-PSS. El JWT se transporta en \emph{cookie} \texttt{HttpOnly} +
\texttt{Secure} + \texttt{SameSite} y las contraseñas se almacenan con un
\emph{hash} seguro~\cite{rfc7519}.

\subsection{A03 --- Injection}

Se comprueba que el sistema rechaza con \texttt{422} las peticiones que no
cumplen la validación de Bean Validation, en lugar de ejecutarlas. La
evidencia A03 muestra la validación de \texttt{POST /api/estudiantes} con
todos los campos obligatorios ausentes, devolviendo \texttt{errores} por
campo en \texttt{application/problem+json}. Las consultas se ejecutan
mediante Spring Data JPA parametrizado y los procedimientos almacenados se
invocan con \texttt{@Procedure}; no se concatena SQL.

\subsection{A04 --- Insecure Design}

El diseño inseguro se previene en la capa de diseño, no solo con \emph{fixes}
puntuales. La evidencia A04 comprueba tres reglas de negocio:

\begin{itemize}[leftmargin=*]
  \item \textbf{Política de contraseñas} aplicada en el servidor:
    \texttt{POST /api/auth/registro} con contraseña corta devuelve
    \texttt{422}.
  \item \textbf{Catálogos obligatorios}: \texttt{POST /api/estudiantes} con
    una categoría inexistente devuelve \texttt{404} (la regla de negocio se
    valida contra el catálogo real, no solo el tipo).
  \item \textbf{Credenciales duplicadas} rechazadas: el registro con un
    \texttt{username} ya usado devuelve \texttt{409}.
\end{itemize}

El modelo de amenazas y las decisiones de diseño se documentan en la
sección~\ref{sec:reflexion}.

\subsection{A05 --- Security Misconfiguration}

Toda respuesta del backend incluye las cabeceras de seguridad:

\begin{lstlisting}[caption={Cabeceras de seguridad (evidencia A05)}]
X-Content-Type-Options: nosniff
X-Frame-Options: DENY
Content-Security-Policy: default-src 'self'; frame-ancestors 'none'
Cache-Control: no-cache, no-store, max-age=0, must-revalidate
\end{lstlisting}

\subsection{A06 --- Vulnerable and Outdated Components}

El inventario de dependencias se verifica con sus versiones y se documenta en
\ruta{docs/mediciones/sec/a06-componentes.txt}: Spring Boot 3.2.5, jjwt
0.12.5, springdoc-openapi 2.3.0, Flyway 10.12.0, Angular 21.2.x y
TypeScript 5.9.2. Las imágenes de PostgreSQL y Redis están fijadas por digest
SHA-256 en \ruta{docker-compose.yml}, de modo que dos construcciones del mismo
\emph{commit} usan exactamente los mismos binarios.

\subsection{A07 --- Identification and Authentication Failures}

La evidencia A07 comprueba el bloqueo por fuerza bruta: los intentos 1--5 de
login con credenciales incorrectas devuelven \texttt{401} y el sexto devuelve
\texttt{429 Too Many Requests}, dejando la IP bloqueada durante 15 minutos.
El controlador de intentos se lleva en Redis por IP de origen
(\texttt{login\_attempts:<ip>}), no por usuario.

\subsection{A09 --- Security Logging and Monitoring Failures}

El registro de auditoría escribe eventos de autenticación con marca de tiempo,
IP de origen y sujeto (evidencia A09): \texttt{AUTH\_LOGIN\_OK} en nivel
\texttt{INFO} y \texttt{AUTH\_LOGIN\_FAIL} en nivel \texttt{WARN}.

\subsection{XSS --- Cross-Site Scripting}

La protección ante XSS se verifica por tres vías (evidencia
\ruta{xss.txt}): (1) la cabecera \texttt{Content-Security-Policy:
default-src 'self'; frame-ancestors 'none'} está presente; (2) un payload
\texttt{<script>} enviado en \texttt{POST /api/categorias} se \emph{almacena}
devuelto con \texttt{201}, pero el servidor lo sirve siempre como
\texttt{application/json}, que el navegador no interpreta como HTML; y (3)
Angular lo escapa por defecto en el SPA al renderizarlo. La evidencia
comprueba que el \emph{payload} viaja como dato escapable, no como \emph{HTML}
ejecutable.

\subsection{Resultado de la auditoría}

\begin{center}
\begin{tabular}{@{}lll@{}}
\toprule
\textbf{Control} & \textbf{Comportamiento verificado} & \textbf{Estado} \\
\midrule
A01 Control de acceso & \texttt{403} en recursos administrativos & \ok \\
A02 Criptografía & TLS 1.3, cookie HttpOnly & \ok \\
A03 Inyección & Validación server-side, JPA parametrizado & \ok \\
A04 Diseño inseguro & Reglas de negocio: 422/404/409 & \ok \\
A05 Configuración & Cabeceras de seguridad presentes & \ok \\
A06 Componentes & Inventario con versiones, digests fijos & \ok \\
A07 Autenticación & 6.º intento $\to$ \texttt{429} (15\,min) & \ok \\
A09 Registro & \texttt{AUTH\_LOGIN\_OK}/\texttt{AUTH\_LOGIN\_FAIL} & \ok \\
XSS & CSP + \texttt{application/json} + escape Angular & \ok \\
\bottomrule
\end{tabular}
\end{center}

\section{Migraciones Flyway y reproducibilidad (Bloque 4)}

Las migraciones versionadas V1--V7 se alinearon con las entidades JPA y con
el esquema consolidado \ruta{db/schema.sql}:

\begin{itemize}[leftmargin=*]
  \item \texttt{V1}: esquema \texttt{academico} + tablas de personas y
    usuarios.
  \item \texttt{V2}: catálogos y \texttt{academico.estudiantes}.
  \item \texttt{V3}: entrenadores con \texttt{id\_usuario}, experiencia y
    certificación, con FK a \texttt{academico}.
  \item \texttt{V4}: claves foráneas a \texttt{academico}.
  \item \texttt{V5}: funciones \texttt{academico} con \texttt{id\_categoria}
    INT.
  \item \texttt{V6}: procedimientos almacenados.
  \item \texttt{V7}: \emph{seed} idempotente (\texttt{admin}/\texttt{Admin2026!}).
\end{itemize}

La verificación se hizo con bases de datos limpias aplicando V1--V7 con
Flyway y arrancando la aplicación con \texttt{ddl-auto=validate}, confirmando
que el esquema migrado coincide con las entidades. El arranque completo del
sistema es reproducible con \texttt{make up}~\cite{flyway}.

\section{Pruebas de carga (Bloque 5.3)}
\label{sec:carga}
\porintegrante{Integrante B. Sección reservada.}

Esta sección documenta las pruebas de carga sobre el \emph{endpoint}
principal de la API con k6~\cite{k6docs}: 3 corridas independientes de 50
usuarios virtuales durante 30 segundos, con análisis de percentiles e
intervalo de confianza al 95\,\%.

\begin{center}
\begin{tabular}{@{}llllll@{}}
\toprule
\textbf{Corrida} & \textbf{P50} & \textbf{P95} & \textbf{P99} & \textbf{req/s} & \textbf{\% error} \\
\midrule
Corrida 1 & [\_\_\_\_] & [\_\_\_\_] & [\_\_\_\_] & [\_\_\_\_] & [\_\_\_\_] \\
Corrida 2 & [\_\_\_\_] & [\_\_\_\_] & [\_\_\_\_] & [\_\_\_\_] & [\_\_\_\_] \\
Corrida 3 & [\_\_\_\_] & [\_\_\_\_] & [\_\_\_\_] & [\_\_\_\_] & [\_\_\_\_] \\
\midrule
Media & [\_\_\_\_] & [\_\_\_\_] & [\_\_\_\_] & [\_\_\_\_] & [\_\_\_\_] \\
IC 95\,\% & [\_\_\_\_] & [\_\_\_\_] & [\_\_\_\_] & [\_\_\_\_] & [\_\_\_\_] \\
\bottomrule
\end{tabular}
\end{center}

\textbf{Umbral objetivo:} P95 menor a 200\,ms con caché caliente y menor a
500\,ms con caché frío (ISO/IEC 25010~\cite{iso25010}).

\porintegrante{Integrante B: completar la tabla con los valores de las
corridas \texttt{k6-run1..3.json}, el análisis del \emph{script}
\ruta{scripts/perf-analysis.py} y la conclusión frente al umbral.}

% ============================================================
\section{Reflexión y modelo de amenazas (Bloque 5.4)}
\label{sec:reflexion}

\subsection{Modelo de amenazas y decisiones de diseño}

El análisis de diseño inseguro (A04) no es una lista de \emph{fixes}, sino un
modelo de amenazas previo a la implementación. Para SGED, cuyo dato más
sensible es la información personal de menores de edad (cédula, fecha de
nacimiento, asistencia y evaluaciones), se decidió:

\begin{itemize}[leftmargin=*]
  \item \textbf{Validación de negocio en el servidor}, no solo en el
    \emph{frontend}: el cliente Angular valida para dar retroalimentación
    inmediata, pero el backend rechaza cualquier violación con
    \texttt{422}/\texttt{404}/\texttt{409} (evidencia A04).
  \item \textbf{Autenticación basada en \emph{cookie} HttpOnly} en lugar de
    JWT en \texttt{localStorage}: el token no es legible por \emph{scripts},
    reduciendo la superficie de exfiltración XSS.
  \item \textbf{Política de contraseñas obligatoria} y bloqueo por fuerza
    bruta con Redis: el límite se aplica por IP para impedir el barrido de
    credenciales sin penalizar a un usuario legítimo.
  \item \textbf{Content-Type siempre \texttt{application/json} + CSP
    \texttt{default-src 'self'}}: un \emph{payload} almacenado no puede
    convertirse en HTML ejecutable; el escape de Angular es la segunda línea
    de defensa.
\end{itemize}

\subsection{Lección transversal: una medición que no se ejecuta contra el
sistema real no es evidencia}

Durante la práctica se materializó tres veces el mismo error de medición:
\emph{un indicador daba verde sobre un sistema que no hacía lo que el
indicador afirmaba}. La primera vez, en la Tercera Entrega, la cobertura de
JaCoCo se medía sobre \texttt{target/} con clases previas a la
reestructuración e incluía paquetes inexistentes. La segunda, la auditoría
A01 devolvía \texttt{401} en todo porque el propio control A07 había bloqueado
la IP de la máquina de auditoría: ese resultado parece correcto y no lo es,
porque probaba autenticación, no autorización. La tercera, el registro de
usuarios estaba roto contra la base de datos mientras sus pruebas unitarias
pasaban, porque usan dobles de prueba y la restricción la aplica PostgreSQL.

La corrección estructural fue la misma en los tres casos: la medición debe
ejecutarse contra el sistema real, desde estado limpio, y comprobar las dos
direcciones (el rechazo seguro \emph{y} el uso legítimo). La auditoría de
esta práctica incorpora ambas propiedades: limpia su propio contador de
intentos al inicio y verifica por cada control que el recurso legítimo sigue
devolviendo \texttt{200}. Esa es la diferencia entre una evidencia que
acompaña a un resultado y una evidencia que lo produce.

\subsection{Contribución de la práctica}

El aporte de la práctica no fue solo añadir documentación y controles, sino
\emph{verificar con evidencia reproducible} lo que el sistema ya afirmaba
tener, y corregir en su origen los desajustes encontrados: las migraciones
Flyway no estaban alineadas con el esquema y hoy lo están; el inventario de
componentes A06 carecía de versiones y hoy las declara; y la auditoría OWASP
ampliada cubre A04, A06, A09 y XSS con la misma disciplina que los controles
originales.

\section{Autoría y honestidad académica}
\label{sec:autorias}

La autoría es verificable en el historial de \texttt{git} de la rama
\texttt{feature/u4-docs-seguridad} (\texttt{git log --author}). Los bloques
de la práctica se asignaron así:

\begin{center}
\begin{tabular}{@{}ll@{}}
\toprule
\textbf{Integrante} & \textbf{Bloques} \\
\midrule
Velez Lopez Ricardo Elias & Bloques 1--4, auditoría OWASP (sección~\ref{sec:seguridad}), informe \\
Integrante A & Bloque 5.1: API externa (sección~\ref{sec:apiexterna}) \\
Integrante B & Bloque 5.3: pruebas de carga (sección~\ref{sec:carga}) \\
\bottomrule
\end{tabular}
\end{center}

\porintegrante{completar nombres, correos institucionales y enlaces a los
\emph{commits} de los integrantes A y B.}

% ============================================================
\section{Conclusiones}

\begin{enumerate}[leftmargin=*]
  \item La API REST de SGED quedó documentada y consultable en
    \texttt{/api/docs} (Swagger UI) y \texttt{/api/api-docs} (OpenAPI 3.0),
    con anotaciones \texttt{@Tag/@Operation/@Schema} en controladores y DTOs.
    La documentación es consumible por herramientas, no un PDF suelto sin
    fuente verificable (Bloques 1--2).

  \item La auditoría de seguridad comprueba 9 controles (A01--A07, A09 y XSS)
    con peticiones \texttt{curl} reproducibles guardadas como evidencia
    cruda. Todos los comportamientos verificados son los esperados, y la
    auditoría comprueba la doble dirección: el rechazo seguro y el uso
    legítimo intacto.

  \item El sistema incorpora defensas de diseño, no solo de código: bloqueo
    por fuerza bruta en Redis, validación de negocio en el servidor, cookie
    HttpOnly y CSP con \texttt{default-src 'self'}. La protección ante XSS se
    evidencia por tres vías independientes.

  \item Las migraciones Flyway V1--V7 quedaron alineadas con las entidades y
    el esquema consolidado, verificadas contra bases de datos limpias con
    \texttt{ddl-auto=validate}, y el arranque completo es reproducible con un
    solo comando (\texttt{make up}).

  \item El uso de herramientas de inteligencia artificial en el desarrollo se
    declara de forma explícita y transparente en la
    sección~\ref{sec:ia}; no sustituye la comprensión ni la autoría, y se
    registra con la misma disciplina que el resto de la práctica.
\end{enumerate}

\section{Trabajo futuro}

\begin{itemize}[leftmargin=*]
  \item Completar los bloques 5.1 (API externa) y 5.3 (pruebas de carga) con
    los resultados de los integrantes A y B, y recompilar este informe.
  \item Incorporar el escaneo OWASP Dependency-Check al CI cuando la base NVD
    sea accesible en el entorno de integración.
  \item Reemplazar el certificado TLS autofirmado por uno de confianza antes
    de cualquier despliegue no académico.
  \item Asignar el DOI de Zenodo a la etiqueta \texttt{v1.0.0-rc} y publicar
    el video de demostración.
  \item Extender la cobertura de pruebas a los módulos nuevos del dominio
    deportivo expuestos en la Entrega Final.
\end{itemize}

\section{Declaración de uso de inteligencia artificial}
\label{sec:ia}

De acuerdo con la guía de transparencia del SWEBOK v4.0~\cite{swebok2024} y
con la política de honestidad académica de la asignatura, se declara
explícitamente que en el desarrollo de esta práctica se utilizaron
herramientas de asistencia por inteligencia artificial. El alcance y los
límites de ese uso se precisan para que la evaluación sea justa:

\begin{itemize}[leftmargin=*]
  \item \textbf{Lo que la IA asistió}: redacción de la estructura de este
    informe en \LaTeX, generación de esqueletos de \emph{scripts} de
    auditoría, verificación de la sintaxis de las migraciones Flyway y ayuda
    en la depuración de mensajes de error de la cadena de herramientas
    (Maven, Flyway, Docker Compose).
  \item \textbf{Lo que no se delegó a la IA}: el diseño del sistema, el
    modelado de amenazas, la escritura de la lógica de negocio, la
    interpretación de los resultados de las mediciones, las decisiones de
    arquitectura y la redacción final de las conclusiones. Cada resultado se
    verificó ejecutando el comando correspondiente contra el sistema real.
  \item \textbf{Traza}: los \emph{commits} de la rama
    \texttt{feature/u4-docs-seguridad} permiten al docente auditar el trabajo
    de cada integrante; la evidencia cruda en \ruta{docs/mediciones/} permite
    reproducir cada cifra citada.
\end{itemize}

% ============================================================
\appendix
\section{Anexo: Taxonomía de Bloom aplicada a la práctica}

Siguiendo la taxonomía revisada de Anderson y Krathwohl~\cite{anderson2001}
(original de Bloom~\cite{bloom1956}), la práctica se ubica principalmente en
los niveles \emph{crear} y \emph{evaluar}: los bloques no solo implementan
una API documentada, sino que construyen y verifican evidencia para juzgar si
el sistema cumple los umbrales exigidos. La tabla relaciona cada nivel con
las actividades concretas de la práctica.

\begin{center}
\begin{tabular}{@{}p{2.6cm}p{6.4cm}p{6.4cm}@{}}
\toprule
\textbf{Nivel} & \textbf{Actividad en la práctica} & \textbf{Artefacto} \\
\midrule
Recordar & Conceptos REST, OWASP, TLS, caché & Secciones 1--3 \\
Comprender & Flujo MVC Spring Boot y JWT en cookie & \ruta{docs/diagramas/} \\
Aplicar & Configurar Swagger, anotar DTOs, escribir migraciones & Bloques 1--4 \\
Analizar & Auditar controles OWASP y descomponer fallos de medición & \ruta{docs/mediciones/sec/} \\
Evaluar & Juzgar cada control frente a umbrales exigidos & Tabla de la sección~\ref{sec:seguridad} \\
Crear & Construir la evidencia reproducible y este informe & \ruta{Makefile}, \ruta{main.tex} \\
\bottomrule
\end{tabular}
\end{center}

\newpage
\bibliographystyle{plain}
\bibliography{referencias}

\end{document}
